\section{Khoảng trống, phát biểu ra}
\label{sec:ch18-khoang-trong-phat-bieu-ra}

\index{khoảng trống nghiên cứu!phát biểu ra|(}%
Ta thấy lĩnh vực đã làm đúng phần việc không đòi trả giá. Mục
này phát biểu khoảng trống thành một câu, nói một đề nghị đang có là bỏ hẳn câu
hỏi, rồi nói những điều cuốn sách không khẳng định. Đây là mục cuối của chương
cuối, nên nó không giao gì cho ai.

\index{khoảng trống nghiên cứu!câu phát biểu}%
Phát biểu thì như sau. Các chỉ số độ trung thực của attribution đặc trưng chưa
bao giờ được đem đối chiếu với một tiêu chí độc lập về việc mô hình thật sự dựa
vào cái gì. Điều đó không phải vì phép đối chiếu ấy không nghĩ ra được: bài siêu
đánh giá nghĩ ra nó và tuyên bố nó không làm được vì thiếu ground truth, còn bài
neo dựng lấy ground truth từ thiết kế nhiệm vụ và chạy nó ở một loại lời giải
thích khác. Cũng không phải vì thiếu số liệu: một bảng in năm 2021 đã đặt bốn chỉ
số thay thế cạnh một cột có tham chiếu và cạnh một mốc ngẫu nhiên, rồi đọc kết
quả thành một câu về các phương pháp. Khoảng trống nằm ở chỗ hai nửa của công
việc chưa gặp nhau: bên có cách dựng ground truth thì dùng nó ở chỗ khác, và bên
cần nó thì đã kết luận rằng không có.

\index{khoảng trống nghiên cứu!cái vòng của ba lối thoát}%
Đọc cùng với mục~\ref{sec:ch16-ba-cai-gia} thì hình dạng ấy quen. Ở đó, ba lối
thoát không xếp thành ba cửa cho một người chọn mà xếp thành một vòng, trong đó
mỗi bên chỉ sang bên khác: người đề xuất đường thứ nhất chỉ sang đường thứ ba và
nói không đi được, người vẽ bản đồ đường thứ ba nói chỗ thiếu nằm ở cả ba nhánh,
người đếm đường thứ hai nói đường ấy đi được nhưng gần như không ai đi. Chương
này thêm một vòng nữa ở tầng dưới, giữa hai bên vừa nói. Cái chung của hai vòng
là không bên nào sai và cũng không bên nào đủ.

\index{khả năng giám sát!một nhánh đề nghị thôi hỏi}%
Có một nhánh không đứng trong vòng ấy, và nó đề nghị bước ra khỏi câu hỏi.
Mục~\ref{sec:ch17-bai-neo-tro-ve} đã thuật lại qua bài neo: có những tác giả cho
rằng định nghĩa cũ không dùng được trên thực tế, nên đề nghị bỏ độ trung thực và
thay bằng \tn{khả năng giám sát}{monitorability}, tức việc tự động dự đoán các
tính chất trong hành vi của mô hình từ chính đầu ra của nó. Theo lời bài neo, các
công trình theo hướng ấy lập luận rằng độ trung thực khó đo được trên thực tế nên
họ kiểm thứ khác~\autocite{p21metrics}.

Sách không đọc hai bài đứng sau đề nghị ấy, và không đọc là một quyết định chứ
không phải một chỗ bỏ sót. Quy tắc mà mục~\ref{sec:ch18-ba-bai-le-ra-da-dong} vừa
dùng ba lần là sách mở một bài ngoài danh mục khi chương phải \emph{xử} một khẳng
định về bài ấy. Ở đây chương không xử: nó ghi nhận rằng đề nghị tồn tại. Nên đề
nghị ấy được thuật lại và không được cân, và điều sách nói về nó chỉ là điều bài
neo nói.

Nhưng việc nó tồn tại là một dữ kiện về hình dạng của khoảng trống, và nó là dữ
kiện cuối cùng mà chương này thêm vào. Mười bảy chương đã gặp những bài đo sai,
những bài không đo, những bài đo thứ khác, và một bài đo được. Đây là nhánh đầu
tiên đề nghị thôi hỏi. Một câu hỏi mà một nhánh của lĩnh vực đề nghị bỏ đi thì
đang ở tình trạng khác với một câu hỏi chưa ai kịp trả lời, và phân biệt hai tình
trạng ấy là việc của người đọc chứ không của cuốn sách.

\index{khoảng trống nghiên cứu!ba điều sách không khẳng định}%
Còn ba điều cuốn sách không khẳng định, và phải nói ra vì mười bảy chương phê
phán dễ bị đọc thành một khẳng định mạnh hơn khẳng định thật.

Thứ nhất, sách không nói các lời giải thích vô ích.
Mục~\ref{sec:ch01-luan-de-cua-sach} đã đặt luận đề hẹp hơn thế ngay từ đầu, và
chương này không mở rộng nó. Một bản đồ nổi bật vẫn chỉ ra chỗ để nhìn, và
chương~\ref{ch:con-nguoi-vang-mat} ghi rằng ở chỗ có người ngồi đọc thì có những
tuyên bố đã được xác nhận, dù không phải tuyên bố về độ trung thực.

Thứ hai, và đây là chỗ dễ trượt nhất, sách không nói các chỉ số \emph{sai}. Chưa
kiểm định không phải là sai. Một dụng cụ chưa được đối chiếu có thể đang đo đúng
thứ ghi trên mặt đồng hồ; điều không có là lý do để tin như vậy. Bài neo tìm được
kết quả gần mức ngẫu nhiên cho tám chỉ số ở phía chuỗi suy luận, và
mục~\ref{sec:ch09-doi-tuong-khac-nhau} đã giữ ranh giới ấy từ đó tới đây để kết
quả ấy không bị mang sang phía này. Nó vẫn không được mang sang. Cái sách khẳng
định là các con số ấy chưa được cấp tư cách bằng chứng, không phải là chúng sai.

Thứ ba, sách không nói khoảng trống này là lỗi của ai. Ba trong chín bài của sổ
không nêu chỗ thiếu vì đối tượng của chúng ở chỗ khác, và điều đó hợp lý. Bài
siêu đánh giá không chạy phép đối chiếu vì nó kết luận rằng không có ground truth
để chạy, và kết luận ấy đúng trong phạm vi nó xét. Bài chuẩn đọc bảng của mình
theo hàng vì bảng ấy được dựng để trả lời câu hỏi theo hàng. Xét từng bước một
thì không bước nào sai. Khoảng trống là thứ hiện ra khi xếp mười bảy chương cạnh
nhau, và
đó là lý do một cuốn sách nhìn thấy nó còn từng bài thì không.

\index{khoảng trống nghiên cứu!phát biểu ra|)}%
Cuốn sách bắt đầu bằng một phân biệt, giữa một lời giải thích nghe hợp lý và một
lời giải thích bám vào cơ chế, và bằng một đòi hỏi, rằng thước đo cho vế thứ hai
phải được kiểm chứng như một dụng cụ đo. Mười bảy chương sau, đòi hỏi ấy vẫn
chưa được đáp ứng ở phía mà nó được đặt ra, đã được đáp ứng hai lần ở hai phía
mà vế phải có sẵn, và có một nhánh đề nghị rút nó lại. Bảng~\ref{tab:ch18-nam-dieu-kien}
là những gì sách gom được về hình dạng của cái sẽ đáp ứng nó, và
mục~\ref{sec:ch18-phep-do-chua-ai-chay} là những gì cần làm để bắt đầu, xếp theo
giá. Cả hai là công việc chứ không phải kết luận, và một cuốn sách về độ trung
thực khép lại đúng ở đó thì trung thực với chính nó.
